problem:  
Let $(M^{2n+1}, \xi, \eta, \Phi, g)$ be a compact, (possibly irregular) Sasaki–Einstein manifold with Reeb field $\xi$. The volume of $M$ and the Reeb vector determine the Hilbert series
\[
H_\xi(t) = \sum_{m\ge 0} \dim \mathcal{H}_m\, e^{-tm},
\]
where $\mathcal{H}_m$ is the space of holomorphic functions on the Kähler cone $C(M)$ homogeneous of weight $m$ under the holomorphic flow generated by $\xi$.  
For quasi-regular $\xi$, $H_\xi(t)$ is the generating function of holomorphic sections of powers of an ample line bundle on the Fano orbifold base; in the irregular case, $H_\xi$ still converges for $\Re(t)>0$ and can be viewed as a holomorphic continuation encoding the “spectrum of Reeb weights.”  

It is known that $H_\xi(t)$ admits an asymptotic expansion near $t=0$ of the form
\[
H_\xi(t) \sim t^{-(n+1)}\big(A_0(\xi) + A_1(\xi)t + A_2(\xi)t^2 + \dots \big),
\]
with $A_0(\xi)$ proportional to the Sasaki volume $\mathrm{Vol}(M,\xi)$. For quasi-regular cases, all $A_i(\xi)$ are computable from orbifold characteristic classes.  

**Open Problem (Irregular Sasaki–Einstein Hilbert asymptotics under Reeb variation):**  
For a smooth family of Reeb fields $\xi_t$ within the Reeb cone of a fixed CR structure $(D,J)$, assume that $\xi_0$ corresponds to a Sasaki–Einstein structure.  
Is the Taylor expansion of each coefficient $A_k(\xi_t)$ in $t$ well-defined and smooth as a function of $\xi_t$ even when $\xi_t$ passes through *irrational* Reeb directions (i.e. when the Reeb flow is non-closed)?  
Equivalently, can the correspondence $\xi \mapsto (A_0(\xi),A_1(\xi),\dots)$ be extended to a real-analytic functional on the entire Reeb cone, in such a way that it agrees with the orbifold expansion in the rational directions?  

If such regularity fails, what analytic singularities (e.g. polylogarithmic or fractal-type) arise in the asymptotics of $H_\xi(t)$ as $\xi$ approaches an irrational Reeb vector?  
How do these singularities reflect the geometry of the underlying Sasaki–Einstein metric and its toric cone, and can they be related to phase transitions in the dual AdS/CFT central charge as a function of $\xi$?

Why is it a "good" problem:  
This problem targets the subtle frontier where smooth Sasakian geometry meets number-theoretic and analytic complexity. The Hilbert series $H_\xi(t)$ is a central invariant closely tied to both the a-maximization functional in AdS/CFT and the Reeb-volume functional in Sasaki–Einstein geometry. While the coefficients $A_k(\xi)$ are understood in rational (quasi-regular) cases through equivariant Riemann–Roch or index theory, the irregular (irrational) regime lies beyond current analytic control.  

Asking whether these asymptotic coefficients extend smoothly or develop singularities probes the deepest interface between differential geometry, spectral analysis, and arithmetic regularity of torus actions—bridging the continuous geometry of the Reeb cone and the discrete structure of rational directions.  

The problem is simple to state (analytic continuation of asymptotic coefficients in $\xi$) yet profoundly difficult: resolving it would require new microlocal or noncommutative-geometric tools to handle irregular Hamiltonian flows and could illuminate how analytic and arithmetic structure intertwine in both Sasaki–Einstein moduli and dual field-theoretic anomaly functions. It thus captures in one apparently modest question an exceptional unification of geometry, analysis, and physics.